\section{Application to land-cover imagery}\label{sec:application}

We now turn to real scenes, where the question is no longer whether a method
recovers a known $K$ but whether it behaves sensibly on data with genuine
spatial autocorrelation, smooth illumination gradients, and spectrally
overlapping classes. We use three benchmarks spanning two sensing modalities
and report two experiments: a direct test of the specificity property of
Section~\ref{subsec:prop} on regions known to contain a single class, and a
full-scene estimate of the number of classes.

\subsection{Data and preprocessing}\label{subsec:data}

Two airborne hyperspectral scenes with per-pixel ground truth are used: the
\emph{Indian Pines} ($145\times145$, $200$ bands after removal of
water-absorption bands, $16$ classes; \citealp{Baumgardner2015}) and the
\emph{Salinas} scene ($512\times217$, $204$ bands, $16$ classes), both standard
in the classification literature. Each cube is standardised per band and
reduced by principal components to ten dimensions, retaining $95.5\%$ and
$99.6\%$ of the variance respectively. As a multispectral counterpart, the
common satellite setting, we use a $246\times237$ Sentinel-2 Level-2A scene
over the Po Valley with ten surface-reflectance bands (B02--B12), labelled by
the ESA WorldCover product \citep{Zanaga2022}; four classes exceed $1\%$ of the
scene (crop $83\%$, tree $9\%$, grass $5\%$, built $4\%$), so its nominal
number of classes is four. Bands are standardised; no dimensionality reduction
is needed.

\subsection{Specificity on single-class regions}\label{subsec:realspec}

A region covered by a single land-cover class is the real-data analogue of the
null and trend worlds of Section~\ref{sec:simulation}: its pixels are
autocorrelated and spectrally variable (soil, moisture and illumination vary
within a crop field), yet the true number of classes is one. We extract all
$16\times16$ windows in which a single ground-truth class covers at least
$90\%$ of the pixels ($67$ windows across Indian Pines and Salinas, $60$ in the
Sentinel-2 scene, all of the dominant crop class), and record the number of
clusters each method reports. Table~\ref{tab:realspec} summarises the outcome;
the proposed-family ablations of Section~\ref{subsec:methods} all behave like
ESS-Dip here (specificity $\ge0.97$) and are omitted for brevity.

\begin{table}[t]
\centering
\caption{Number of clusters reported on single-class windows (true $K=1$).
A method calibrated for spatial autocorrelation should return $\hat K=1$.}
\label{tab:realspec}
\begin{tabular}{l cc cc}
\toprule
 & \multicolumn{2}{c}{Hyperspectral ($67$ win.)}
 & \multicolumn{2}{c}{Sentinel-2 ($60$ win.)}\\
\cmidrule(lr){2-3}\cmidrule(lr){4-5}
Method & $\hat K{=}1$ rate & mean $\hat K$ & $\hat K{=}1$ rate & mean $\hat K$\\
\midrule
Gap statistic       & 0.64 & 1.79 & 0.00 & 6.97\\
Silhouette          & 0.00 & 2.27 & 0.00 & 2.35\\
Calinski--Harabasz  & 0.00 & 2.70 & 0.00 & 3.20\\
Davies--Bouldin     & 0.00 & 3.37 & 0.00 & 2.62\\
\midrule
ESS-Dip             & \textbf{0.99} & \textbf{1.01} & \textbf{1.00} & \textbf{1.00}\\
\bottomrule
\end{tabular}
\end{table}

The simulation result reproduces on real imagery. The silhouette,
Calinski--Harabasz and Davies--Bouldin indices split a single real land-cover
class into two or more sub-clusters in \emph{every} window of both modalities;
the gap statistic does so in all but the cleanest hyperspectral windows
($\hat K{=}1$ rate $0.64$) and, on the multispectral crop fields, fragments a
single class into nearly seven on average. ESS-Dip identifies the region as a
single class in $99\%$ of the hyperspectral and $100\%$ of the multispectral
windows. The within-field spectral variability that the standard indices read
as structure is precisely what the effective-sample-size calibration discounts.
The result is not an artefact of the preprocessing: on the Indian Pines windows
ESS-Dip returns a single class in every window across principal-component
truncations from five to thirty components, window sides from $12$ to $24$
pixels, and purity thresholds from $0.85$ to $0.95$ (for those combinations that
yield windows). The recursive Gaussian-simulation estimator,
even with the multiple-testing correction of
Section~\ref{subsec:geostat-ref}, is less reliable here: it returns a single
class in $0.82$ of the hyperspectral and $0.98$ of the multispectral windows,
below ESS-Dip, because its per-node plug-in variogram degrades on real,
non-stationary support, the residual effect that the correction does not
remove.

\subsection{Full-scene estimation}\label{subsec:realfull}

Applied to an entire labelled scene, the methods estimate the number of
land-cover classes present (Table~\ref{tab:realfull}). No criterion approaches
the nominal sixteen classes of the hyperspectral scenes: the labels include
many small, spectrally similar agricultural classes that are not separable
without supervision, and the best standard index reaches nine (the gap on
Salinas) while the others return two to six. ESS-Dip is deliberately
conservative, reporting two classes on Indian Pines and five on Salinas. On the
crop-dominated Sentinel-2 scene, where a single class covers $83\%$ of the
pixels, ESS-Dip returns that one dominant class, whereas the gap reports five.

\begin{table}[t]
\centering
\caption{Estimated number of classes on full scenes; the nominal number of
labelled classes is given in parentheses. The $\ge20$ entry marks the recursive
simulation estimator reaching its recursion cap without terminating.}
\label{tab:realfull}
\begin{tabular}{l ccc}
\toprule
Method & Indian Pines (16) & Salinas (16) & Sentinel-2 (4)\\
\midrule
Gap statistic       & 3 & 9 & 5\\
Silhouette          & 2 & 2 & 2\\
Calinski--Harabasz  & 2 & 6 & 3\\
Davies--Bouldin     & 2 & 6 & 2\\
ESS-Dip             & 2 & 5 & 1\\
SGS-Dip (corrected) & 2 & $\ge20$ & 6\\
\bottomrule
\end{tabular}
\end{table}

These full-scene figures should be read against the precision--recall
trade-off of Section~\ref{subsec:simresults}. Exhaustive recovery of many
overlapping classes from a single scene is beyond every criterion considered,
and the conservatism of ESS-Dip is the price of its specificity: a procedure
that refuses to invent structure will, on a scene dominated by one class or
composed of weakly separated ones, under-report rather than over-report. The
recursive simulation estimator errs in the opposite direction, over-detecting
to the imposed cap on Salinas and to six classes on the crop-dominated
Sentinel-2 scene even after the multiple-testing correction, the plug-in
variogram on the scene's irregular, non-stationary support compounding over a
full-scene recursion. We
regard the method's value as residing in that precision (it answers
reliably whether genuine class structure is present beyond spatial
autocorrelation) rather than in maximal class enumeration, and we return to
this distinction, and to ways of relaxing it, in Section~\ref{sec:discussion}.

The number of classes is only half the picture; the quality of the induced
partition completes it. Table~\ref{tab:realari} reports the adjusted Rand index
(ARI) and normalised mutual information (NMI) of each labelling against the
ground-truth map. Two things stand out. First, unsupervised recovery of these
partitions is hard for every method: no ARI exceeds $0.52$. Second, on the
multi-class hyperspectral scenes ESS-Dip's partition quality trails the
over-detecting gap statistic (Salinas ARI $0.32$ versus $0.52$), and on the
crop-dominated Sentinel-2 scene, where it returns a single class, its ARI is
zero. This is the partition-quality counterpart of the conservatism already
visible in the class counts: a precision-oriented test that refuses to split
does not, and is not meant to, maximise partition recovery. The numbers make
the trade-off explicit and reinforce the reading of ESS-Dip as a homogeneity
guard rather than a classifier.

\begin{table}[t]
\centering
\caption{Partition quality on the full scenes: adjusted Rand index (ARI) and
normalised mutual information (NMI) against the ground-truth land-cover map,
with the number of classes each method returns. Computed on labelled pixels.}
\label{tab:realari}
\begin{tabular}{l l c cc}
\toprule
Scene & Method & $\hat K$ & ARI & NMI\\
\midrule
Indian Pines & Gap statistic & 3 & 0.21 & 0.36\\
             & Silhouette    & 2 & 0.18 & 0.37\\
             & ESS-Dip       & 2 & 0.18 & 0.37\\
\midrule
Salinas      & Gap statistic & 9 & \textbf{0.52} & \textbf{0.71}\\
             & Silhouette    & 2 & 0.18 & 0.37\\
             & ESS-Dip       & 5 & 0.32 & 0.57\\
\midrule
Sentinel-2   & Gap statistic & 5 & 0.10 & 0.14\\
             & Silhouette    & 2 & 0.36 & 0.22\\
             & ESS-Dip       & 1 & 0.00 & 0.00\\
\bottomrule
\end{tabular}
\end{table}
